\chapter*{Introduction générale}
\addcontentsline{toc}{chapter}{Introduction générale}

L'évolution des technologies de l'information a profondément transformé les modes de gestion dans de nombreux domaines, y compris celui des achats et de la chaîne d'approvisionnement. Aujourd'hui, l'automatisation des processus internes et l'intégration de l'intelligence artificielle sont devenues une nécessité pour les entreprises. Elles permettent d'améliorer la fiabilité des décisions, la transparence et la rapidité de traitement.

Dans ce contexte, un projet a été mis en place pour automatiser et fiabiliser le cycle d'approvisionnement au sein de l'ERP \textit{Microsoft Dynamics 365 Business Central}. Il a pour objectif principal de sécuriser et d'objectiver la décision d'achat, en s'appuyant sur des solutions intelligentes intégrées à l'ERP.

Le présent rapport s'inscrit dans le cadre d'un stage pratique réalisé au sein de la société \textbf{DataWise}. Il décrit les différentes phases de réalisation du projet, de l'analyse des besoins à l'implémentation de la solution.

Ce document est structuré en trois grandes parties :

\begin{itemize}
    \item[---] \textbf{La première partie : Présentation générale.} Elle expose une brève présentation de l'organisme d'accueil, le projet et les missions qui m'ont été confiées durant le stage.
    \item[---] \textbf{La deuxième partie : Conception et réalisation.} Elle décrit les étapes de conception du projet et ses aspects techniques.
    \item[---] \textbf{La troisième partie : Conclusion et perspectives.} Elle met en évidence les acquis du stage et propose des améliorations possibles.
\end{itemize}

Ce rapport reflète à la fois l'approche méthodologique adoptée et les compétences techniques et professionnelles développées tout au long du stage.